\documentclass[12pt, letterpaper]{article}
\usepackage[utf8]{inputenc}
\usepackage{amsmath, amssymb, geometry}
\geometry{margin=1in}

\title{\textbf{Solution Key: Perturbations Workshop}}
\author{Aerospace Engineering Program \\ Universidad de Antioquia}
\date{}

\begin{document}

\maketitle

\section*{Solution 1: The Baseline $J_2$ Effect}
Given: Circular orbit ($e=0$), $a = p = 30,000$ km.
\begin{enumerate}
    \item[a)] \textbf{Mean motion:} \\
    $n = \sqrt{\frac{\mu}{a^3}} = \sqrt{\frac{398600}{30000^3}} = 1.215 \times 10^{-4} \text{ rad/s}$.
    
    \item[b)] \textbf{Equatorial Orbit ($i = 0^\circ$):} \\
    $\cos(0^\circ) = 1$, $\sin(0^\circ) = 0$. \\
    $\dot{\Omega} = -\frac{3}{2} n J_2 \left(\frac{R}{p}\right)^2 \cos i$ \\
    $\dot{\Omega} = -1.5 (1.215 \times 10^{-4})(1.0826 \times 10^{-3}) \left(\frac{6378}{30000}\right)^2 (1) = -8.91 \times 10^{-9} \text{ rad/s}$. \\
    $\dot{\omega} = \frac{3}{2} n J_2 \left(\frac{R}{p}\right)^2 \left(2 - \frac{5}{2}\sin^2 i\right)$ \\
    $\dot{\omega} = 1.5 (1.215 \times 10^{-4})(1.0826 \times 10^{-3}) \left(\frac{6378}{30000}\right)^2 (2 - 0) = 1.78 \times 10^{-8} \text{ rad/s}$.
    
    \item[c)] \textbf{Polar Orbit ($i = 90^\circ$):} \\
    $\cos(90^\circ) = 0$, $\sin(90^\circ) = 1$. \\
    $\dot{\Omega} = 0 \text{ rad/s}$. (Physical significance: The oblateness bulge pulls radially outward on a polar orbit but provides no lateral torque to precess the plane). \\
    $\dot{\omega} = 8.91 \times 10^{-9} \left(2 - 2.5(1)\right) = -4.45 \times 10^{-9} \text{ rad/s}$.
\end{enumerate}

\section*{Solution 2: Designing a Sun-Synchronous Orbit (SSO)}
Given: $h = 800$ km $\implies a = p = 7178$ km.
\begin{enumerate}
    \item[a)] \textbf{Convert Precession Rate:} \\
    $\dot{\Omega}_{req} = 0.9856 \frac{\text{deg}}{\text{day}} \times \frac{\pi \text{ rad}}{180 \text{ deg}} \times \frac{1 \text{ day}}{86400 \text{ s}} = 1.991 \times 10^{-7} \text{ rad/s}$.
    
    \item[b)] \textbf{Calculate Inclination:} \\
    $n = \sqrt{\frac{398600}{7178^3}} = 1.038 \times 10^{-3} \text{ rad/s}$. \\
    Rearranging the nodal precession equation for $\cos i$: \\
    $\cos i = \frac{\dot{\Omega}}{-\frac{3}{2} n J_2 \left(\frac{R}{p}\right)^2}$ \\
    $\cos i = \frac{1.991 \times 10^{-7}}{-1.5 (1.038 \times 10^{-3}) (1.0826 \times 10^{-3}) \left(\frac{6378}{7178}\right)^2} = \frac{1.991 \times 10^{-7}}{-1.330 \times 10^{-6}} = -0.1497$ \\
    $i = \arccos(-0.1497) = 98.61^\circ$.
\end{enumerate}

\section*{Solution 3: Critical Inclination and Frozen Orbits}
\begin{enumerate}
    \item[a)] \textbf{Mathematical Derivation:} \\
    For the argument of perigee to remain fixed, $\dot{\omega} = 0$. \\
    The governing equation is $\dot{\omega} = \frac{3}{2} n J_2 \left(\frac{R}{p}\right)^2 \left(2 - \frac{5}{2}\sin^2 i\right)$. \\
    Since $n, J_2, R$, and $p$ are non-zero for any practical orbit, the term in the parenthesis must equal zero: \\
    $2 - \frac{5}{2}\sin^2 i = 0 \implies \frac{5}{2}\sin^2 i = 2 \implies \sin^2 i = \frac{4}{5} = 0.8$ \\
    $\sin i = \pm \sqrt{0.8} \implies i = \arcsin(\sqrt{0.8})$. \\
    Evaluating this gives the two critical inclinations: $i_1 \approx 63.43^\circ$ and $i_2 \approx 116.57^\circ$.
\end{enumerate}

\section*{Solution 4: Interplanetary Perturbations - Jupiter SSO}
\begin{enumerate}
    \item[a)] \textbf{Jovian SSO Precession Rate:} \\
    $\dot{\Omega}_{req} = \frac{360^\circ}{4332.589 \text{ days}} = 0.08309^\circ/\text{day}$. \\
    Converting to rad/s: $0.08309 \times \left(\frac{\pi}{180}\right) \times \left(\frac{1}{86400}\right) = 1.678 \times 10^{-8} \text{ rad/s}$.
    
    \item[b)] \textbf{Orbit Design:} \\
    Chosen semi-major axis: $a = 2 \times 69911 = 139822$ km. ($e=0 \implies p = 139822$ km). \\
    $n = \sqrt{\frac{1.26686 \times 10^8}{139822^3}} = 2.152 \times 10^{-4} \text{ rad/s}$. \\
    $\cos i = \frac{1.678 \times 10^{-8}}{-1.5 (2.152 \times 10^{-4}) (0.014736) \left(\frac{69911}{139822}\right)^2}$ \\
    $\cos i = \frac{1.678 \times 10^{-8}}{-1.189 \times 10^{-6}} = -0.0141$ \\
    $i = \arccos(-0.0141) = 90.81^\circ$. (Notice that due to Jupiter's large distance from the Sun, making the required precession very small, the SSO is nearly a true polar orbit).
\end{enumerate}

\section*{Solution 5: Orbital Decay via Atmospheric Drag}
\begin{enumerate}
    \item[a)] \textbf{Ballistic Coefficient:} \\
    $B = \frac{C_D A}{m} = \frac{2.0 \times 0.03 \text{ m}^2}{3.0 \text{ kg}} = 0.02 \text{ m}^2/\text{kg}$.
    
    \item[b)] \textbf{Decay Rate in m/s:} \\
    First, find mean motion $n$ in standard SI units: \\
    $n = \sqrt{\frac{3.986 \times 10^{14} \text{ m}^3/\text{s}^2}{(6878000 \text{ m})^3}} = 1.107 \times 10^{-3} \text{ rad/s}$. \\
    $\frac{da}{dt} = -n a^2 \rho B$ \\
    $\frac{da}{dt} = -(1.107 \times 10^{-3})(6878000)^2 (1.0 \times 10^{-12}) (0.02)$ \\
    $\frac{da}{dt} = -0.001047 \text{ m/s}$.
    
    \item[c)] \textbf{Daily Altitude Loss:} \\
    $\Delta a_{\text{day}} = -0.001047 \frac{\text{m}}{\text{s}} \times 86400 \frac{\text{s}}{\text{day}} = -90.46 \text{ m/day}$. \\
    The CubeSat loses approximately 90.5 meters of altitude in its first day.
\end{enumerate}

\end{document}